% !TEX root = ../algebra_02.tex

\section{Свойства матриц перехода между базисами}

Пусть $V$ --- конечномерное линейное пространство над полем $K$.

\begin{Definition}{Координатный столбец}{}
	Для $v\in V$ координатный столбец $[v]_E\in K^n$ определяется равенством
	$v=E[v]_E$, то есть
	\[
	v=e_1x_1+\cdots+e_nx_n
	\quad \Longleftrightarrow \quad
	[v]_E=
	\begin{pmatrix}
		x_1\\
		\vdots\\
		x_n
	\end{pmatrix}.
	\]
\end{Definition}

\begin{Definition}{Матрица линейного отображения}{}
	Пусть $\alpha\colon V\to W$ --- линейное отображение. Его матрицей в базисах $E,F$ называется матрица $[\alpha]_{E,F}\in M_{m\times n}(K)$, столбцы которой равны координатам образов базисных векторов:
	\[
	[\alpha]_{E,F}
	=
	\bigl(
	[\alpha e_1]_F,\ldots,[\alpha e_n]_F
	\bigr).
	\]
	Иначе говоря,
	\[
	\alpha E=F[\alpha]_{E,F}.
	\]
\end{Definition}

\begin{Definition}{Матрица перехода}{}
	Пусть $E$ и $E'$ --- два базиса пространства $V$. Матрицей перехода от базиса $E'$ к базису $E$ называется матрица тождественного оператора $\operatorname{id}_V$ в этих базисах:
	\[
	C_{E'\to E} := [\operatorname{id}_V]_{E',E}.
	\]
\end{Definition}

Матрица перехода обладает следующими ключевыми свойствами, которые вытекают из теоремы о матрице композиции линейных отображений.

\begin{Theorem}{Свойства матриц перехода}{}
	Пусть $E, E', E''$ --- базисы линейного пространства $V$. Тогда:
	\begin{enumerate}
		\item Тождественный переход: $C_{E\to E} = I_n$, где $I_n$ --- единичная матрица.
		\item Транзитивность (композиция переходов): $C_{E''\to E} = C_{E'\to E} C_{E''\to E'}$.
		\item Обратимость: Матрица перехода всегда обратима, причём $(C_{E'\to E})^{-1} = C_{E\to E'}$.
	\end{enumerate}
\end{Theorem}

\begin{proof}
	\begin{enumerate}
		\item По определению, $C_{E\to E} = [\operatorname{id}_V]_{E,E}$. Поскольку $\operatorname{id}_V(e_i) = e_i$, координатные столбцы базисных векторов в своём же базисе --- это столбцы единичной матрицы. Следовательно, $C_{E\to E} = I_n$.
		
		\item Рассмотрим композицию отображений $\operatorname{id}_V \circ \operatorname{id}_V = \operatorname{id}_V$. Применим формулу матрицы композиции (для цепочки $V \xrightarrow{\operatorname{id}_V} V \xrightarrow{\operatorname{id}_V} V$ с базисами $E'', E', E$ соответственно):
		\[
		[\operatorname{id}_V]_{E'',E} = [\operatorname{id}_V]_{E',E} [\operatorname{id}_V]_{E'',E'}.
		\]
		По определению матрицы перехода, это равенство записывается как $C_{E''\to E} = C_{E'\to E} C_{E''\to E'}$.
		
		\item Запишем свойство транзитивности для базисов $E, E', E$:
		\[
		C_{E\to E} = C_{E'\to E} C_{E\to E'}.
		\]
		Так как $C_{E\to E} = I_n$, получаем $C_{E'\to E} C_{E\to E'} = I_n$. Аналогично, поменяв местами $E$ и $E'$, получим $C_{E\to E'} C_{E'\to E} = I_n$. Это означает, что матрица $C_{E'\to E}$ обратима, и её обратная матрица равна $C_{E\to E'}$.
	\end{enumerate}
\end{proof}
